% Options for packages loaded elsewhere
% Options for packages loaded elsewhere
\PassOptionsToPackage{unicode}{hyperref}
\PassOptionsToPackage{hyphens}{url}
\PassOptionsToPackage{dvipsnames,svgnames,x11names}{xcolor}
%
\documentclass[
  english,
  11pt,
]{article}
\usepackage{xcolor}
\usepackage[margin=25mm]{geometry}
\usepackage{amsmath,amssymb}
\setcounter{secnumdepth}{-\maxdimen} % remove section numbering
\usepackage{iftex}
\ifPDFTeX
  \usepackage[T1]{fontenc}
  \usepackage[utf8]{inputenc}
  \usepackage{textcomp} % provide euro and other symbols
\else % if luatex or xetex
  \usepackage{unicode-math} % this also loads fontspec
  \defaultfontfeatures{Scale=MatchLowercase}
  \defaultfontfeatures[\rmfamily]{Ligatures=TeX,Scale=1}
\fi
\usepackage{lmodern}
\ifPDFTeX\else
  % xetex/luatex font selection
\fi
% Use upquote if available, for straight quotes in verbatim environments
\IfFileExists{upquote.sty}{\usepackage{upquote}}{}
\IfFileExists{microtype.sty}{% use microtype if available
  \usepackage[]{microtype}
  \UseMicrotypeSet[protrusion]{basicmath} % disable protrusion for tt fonts
}{}
\makeatletter
\@ifundefined{KOMAClassName}{% if non-KOMA class
  \IfFileExists{parskip.sty}{%
    \usepackage{parskip}
  }{% else
    \setlength{\parindent}{0pt}
    \setlength{\parskip}{6pt plus 2pt minus 1pt}}
}{% if KOMA class
  \KOMAoptions{parskip=half}}
\makeatother
% Make \paragraph and \subparagraph free-standing
\makeatletter
\ifx\paragraph\undefined\else
  \let\oldparagraph\paragraph
  \renewcommand{\paragraph}{
    \@ifstar
      \xxxParagraphStar
      \xxxParagraphNoStar
  }
  \newcommand{\xxxParagraphStar}[1]{\oldparagraph*{#1}\mbox{}}
  \newcommand{\xxxParagraphNoStar}[1]{\oldparagraph{#1}\mbox{}}
\fi
\ifx\subparagraph\undefined\else
  \let\oldsubparagraph\subparagraph
  \renewcommand{\subparagraph}{
    \@ifstar
      \xxxSubParagraphStar
      \xxxSubParagraphNoStar
  }
  \newcommand{\xxxSubParagraphStar}[1]{\oldsubparagraph*{#1}\mbox{}}
  \newcommand{\xxxSubParagraphNoStar}[1]{\oldsubparagraph{#1}\mbox{}}
\fi
\makeatother


\usepackage{longtable,booktabs,array}
\newcounter{none} % for unnumbered tables
\usepackage{calc} % for calculating minipage widths
% Correct order of tables after \paragraph or \subparagraph
\usepackage{etoolbox}
\makeatletter
\patchcmd\longtable{\par}{\if@noskipsec\mbox{}\fi\par}{}{}
\makeatother
% Allow footnotes in longtable head/foot
\IfFileExists{footnotehyper.sty}{\usepackage{footnotehyper}}{\usepackage{footnote}}
\makesavenoteenv{longtable}
\usepackage{graphicx}
\makeatletter
\newsavebox\pandoc@box
\newcommand*\pandocbounded[1]{% scales image to fit in text height/width
  \sbox\pandoc@box{#1}%
  \Gscale@div\@tempa{\textheight}{\dimexpr\ht\pandoc@box+\dp\pandoc@box\relax}%
  \Gscale@div\@tempb{\linewidth}{\wd\pandoc@box}%
  \ifdim\@tempb\p@<\@tempa\p@\let\@tempa\@tempb\fi% select the smaller of both
  \ifdim\@tempa\p@<\p@\scalebox{\@tempa}{\usebox\pandoc@box}%
  \else\usebox{\pandoc@box}%
  \fi%
}
% Set default figure placement to htbp
\def\fps@figure{htbp}
\makeatother



\ifLuaTeX
\usepackage[bidi=basic,shorthands=off]{babel}
\else
\usepackage[bidi=default,shorthands=off]{babel}
\fi
\ifLuaTeX
  \usepackage{selnolig} % disable illegal ligatures
\fi


\setlength{\emergencystretch}{3em} % prevent overfull lines

\providecommand{\tightlist}{%
  \setlength{\itemsep}{0pt}\setlength{\parskip}{0pt}}



 


\makeatletter
\@ifpackageloaded{caption}{}{\usepackage{caption}}
\AtBeginDocument{%
\ifdefined\contentsname
  \renewcommand*\contentsname{Table of contents}
\else
  \newcommand\contentsname{Table of contents}
\fi
\ifdefined\listfigurename
  \renewcommand*\listfigurename{List of Figures}
\else
  \newcommand\listfigurename{List of Figures}
\fi
\ifdefined\listtablename
  \renewcommand*\listtablename{List of Tables}
\else
  \newcommand\listtablename{List of Tables}
\fi
\ifdefined\figurename
  \renewcommand*\figurename{Figure}
\else
  \newcommand\figurename{Figure}
\fi
\ifdefined\tablename
  \renewcommand*\tablename{Table}
\else
  \newcommand\tablename{Table}
\fi
}
\@ifpackageloaded{float}{}{\usepackage{float}}
\floatstyle{ruled}
\@ifundefined{c@chapter}{\newfloat{codelisting}{h}{lop}}{\newfloat{codelisting}{h}{lop}[chapter]}
\floatname{codelisting}{Listing}
\newcommand*\listoflistings{\listof{codelisting}{List of Listings}}
\makeatother
\makeatletter
\makeatother
\makeatletter
\@ifpackageloaded{caption}{}{\usepackage{caption}}
\@ifpackageloaded{subcaption}{}{\usepackage{subcaption}}
\makeatother
\usepackage{bookmark}
\IfFileExists{xurl.sty}{\usepackage{xurl}}{} % add URL line breaks if available
\urlstyle{same}
% fallback for those not using the hyperref driver hyperxmp:
\makeatletter
\@ifundefined{xmpquote}{\newcommand{\xmpquote}[1]{#1}}{}
\makeatother
\hypersetup{
  pdftitle={The Navier--Stokes paper{,} one building block at a time},
  pdflang={en},
  colorlinks=true,
  linkcolor={blue},
  filecolor={Maroon},
  citecolor={Blue},
  urlcolor={Blue},
  pdfcreator={LaTeX via pandoc}}


\title{The Navier--Stokes paper, one building block at a time}
\usepackage{etoolbox}
\makeatletter
\providecommand{\subtitle}[1]{% add subtitle to \maketitle
  \apptocmd{\@title}{\par {\large #1 \par}}{}{}
}
\makeatother
\subtitle{Start here: the question, the plan, and a route through the
explanations}
\author{}
\date{2026-09-12}
\begin{document}
\maketitle


A long mathematical paper can feel like a box of LEGO pieces without the
picture on the lid. You can inspect each piece and still have little
idea what you are trying to build. This guide supplies that picture.

The paper describes a specially designed fluid motion that becomes
faster and more concentrated as a particular time approaches. The
challenge is to make that motion obey the fluid equations while the
external force stays smooth. The many profiles, waves, corrections, and
estimates all serve that goal.

You can follow the explanations here with arithmetic, graphs, areas, and
the idea of a slope. When a more advanced concept appears, we explain
its job before naming it. The calculations and original technical
articles remain available at the end of each installment.

\subsection{The question in ordinary
language}\label{the-question-in-ordinary-language}

The Navier--Stokes equations describe how a viscous fluid moves.
\emph{Viscous} means that neighbouring regions moving at different
speeds exchange momentum, tending to smooth out those differences.

The equations combine two ideas. First, fluid parcels accelerate
according to the pushes they receive: pressure, viscosity, and an
external force. Second, the incompressible model makes each moving
parcel keep its volume. Fluid entering a small region must be balanced
by fluid leaving it.

The paper's claimed result concerns a deliberately chosen \textbf{smooth
external force}, acting on fluid initially at rest. The resulting
velocity is smooth before the designated time, but its largest value
becomes unbounded as that time approaches. The total kinetic energy
stays bounded. This is a statement about the forced mathematical
equations.
\href{https://cdn.openai.com/pdf/32d9f210-8b73-45e0-91bc-82a30aef8a9a/navier-stokes.pdf}{Paper,
Theorem 1.1 and Sections 1--2}

``Blow-up'' means that no finite number remains an upper limit on the
speed. It does not mean that the whole fluid suddenly has infinite
speed. Nor does this guide establish that a real liquid or gas could
reach that limit. We are explaining the construction presented in the
paper, not independently checking every proof.

\subsection{The picture on the box}\label{the-picture-on-the-box}

Here is the whole strategy in one sentence:

\textbf{Design a shrinking vortex, calculate the push that its motion is
missing, supply that push through carefully arranged small-scale motion,
and control the errors so the external force stays smooth.}

\begin{figure}

\centering{

\pandocbounded{\includegraphics[keepaspectratio,alt={Six connected blocks: concentrate a vortex; join core to exterior; identify missing transport; prepare shear and add pulses; repair remaining errors; localize the result and keep the force smooth.}]{figures/reader.png}}

}

\caption{\label{fig-map}The argument runs from a concentrating vortex to
a smooth external force. The repair stage repeatedly addresses errors
introduced by earlier repairs. This is a conceptual map, not a
simulation.}

\end{figure}%

\subsubsection{Block 1: a fast motion in a shrinking
region}\label{block-1-a-fast-motion-in-a-shrinking-region}

A small amount of fluid can move extremely fast without making the total
energy enormous. A simple calculation shows why: energy depends on speed
squared \textbf{and} how much fluid moves at that speed.

The intense part of the proposed vortex shrinks in radius and length.
The radius shrinks slightly faster, making the region increasingly
slender. ``The core'' names this moving region of intense motion; it
does not identify a fixed collection of molecules.

\textbf{Job in the argument:} make growing speed compatible with bounded
energy. Start with \href{../d2_finite_energy/article.html}{speed and
energy}. This is mainly Section 2.1, developed in Sections 3--4.

\subsubsection{Block 2: a core, a joining ring, and an
exterior}\label{block-2-a-core-a-joining-ring-and-an-exterior}

The flow needs a smooth centre. Outside it, the motion must connect to a
simpler outer flow. The ring-shaped joining region is called the
\emph{annulus}. Pressure also has to fit across the regions.

Think of building a bridge: an excellent central span is insufficient if
its connections to the banks do not fit. Here the connections must
satisfy fluid balances as well as match smoothly.

\textbf{Job in the argument:} provide a usable background vortex. Read
\href{../d6_inner_profile/article.html}{the smooth centre} and
\href{../d6_axis_pressure/article.html}{pressure from the surrounding
flow}. These explain parts of Section 4 and Appendices A--B.

\subsubsection{Block 3: find the missing
push}\label{block-3-find-the-missing-push}

Once a candidate velocity and pressure have been chosen, we can
calculate the external force they would require. The paper calls this
the \emph{momentum residual}. It is a force-balance calculation.

The background vortex alone leaves a troublesome imbalance in its
joining ring. Calling that imbalance an external force would not solve
the problem: the force would have the wrong behaviour near the singular
time.

\textbf{Job in the argument:} turn ``this vortex is not yet a solution
with a smooth force'' into a precise repair target. Read
\href{../d3_residual/article.html}{the missing-push calculation}. See
Sections 2.2, 3.2 and 5.

\subsubsection{Block 4: use fine ripples to achieve what a coarse shape
cannot}\label{block-4-use-fine-ripples-to-achieve-what-a-coarse-shape-cannot}

A line can look almost flat from far away while being steep on each tiny
rising or falling segment. This separation between overall shape and
local slope is one route into \textbf{convex integration}.

Appendix C uses a related construction: rapid radial variations improve
the local shear---the difference between neighbouring velocities---while
making only a small change to the overall profiles. Their averages are
chosen carefully, and further adjustments restore the required totals.

\textbf{Job in the argument:} prepare a background whose missing
transport can be supplied by the available waves. Read
\href{../convex-integration/article.html}{convex integration, from
slopes to fluid motion}. Appendix C belongs here in the story, even
though it appears near the back of the paper.

\subsubsection{Block 5: make oscillations do useful
work}\label{block-5-make-oscillations-do-useful-work}

Opposite velocity fluctuations can average to zero while their momentum
transport does not. The reason is simple multiplication: reversing two
signs leaves their product unchanged.

The paper chooses two pulse families whose combined transport supplies
the required leading contribution. Their growth comes from the
background shear. Later, shear squeezes their pattern into finer scales
and viscosity damps them.

\textbf{Job in the argument:} let internal fluid motion supply the
troublesome missing transport. Read
\href{../d4_mean_stress/article.html}{zero average, nonzero transport},
then \href{../d5_shearing_wave/article.html}{why a pulse grows and
fades}. These connect to Sections 6--7 and Appendix C.

\subsubsection{Block 6: repair the repairs, then finish the
construction}\label{block-6-repair-the-repairs-then-finish-the-construction}

Adding waves fixes one imbalance and creates others. Local adjustments
can also change totals that matter outside their own neighbourhood. The
construction repeatedly updates the waves, mean motion, pressure, and
weighted totals.

Finally, the paper combines the corrections and confines the flow to a
bounded region. It must preserve both the increasing velocity and the
smoothness of the required external force.

\textbf{Job in the argument:} turn the mechanism into the claimed
complete solution. Read
\href{../d6_moment_correction/article.html}{repairing weighted totals},
then \href{../completing-the-flow/article.html}{from corrections to a
complete flow}. These explain the roles of Sections 5, 8, 9 and 10.

\subsection{Where every section
belongs}\label{where-every-section-belongs}

The order of a proof follows its dependencies. A first reading can
instead follow the ideas.

{\def\LTcaptype{none} % do not increment counter
\begin{longtable}[]{@{}
  >{\raggedright\arraybackslash}p{(\linewidth - 2\tabcolsep) * \real{0.5000}}
  >{\raggedright\arraybackslash}p{(\linewidth - 2\tabcolsep) * \real{0.5000}}@{}}
\toprule\noalign{}
\begin{minipage}[b]{\linewidth}\raggedright
Paper location
\end{minipage} & \begin{minipage}[b]{\linewidth}\raggedright
What it does, in everyday language
\end{minipage} \\
\midrule\noalign{}
\endhead
\bottomrule\noalign{}
\endlastfoot
Section 1 & States the target and relates it to earlier work. \\
Section 2 & Describes the vortex, the pulses, and the physical
mechanism. \\
Section 3 & Gives the plan of the construction before the full
details. \\
Section 4 & Builds the basic vortex and specifies what the joining ring
must do. \\
Section 5 & Adds finer background corrections for balances omitted at
first. \\
Section 6 & Organizes the oscillations so distinct pulses do not create
unwanted overlap terms. \\
Section 7 & Builds growing and decaying pulses and chooses their
averaged transport. \\
Section 8 & Repairs mean motion, pressure compatibility, and weighted
totals locally. \\
Section 9 & Repeats the repairs, controls new errors, and combines the
stages. \\
Section 10 & Starts from rest, confines the construction, extends the
force smoothly, and completes the argument. \\
Appendix A & Prepares radial schedules, matching totals, and the outer
flow. \\
Appendix B & Builds a smooth inner profile and connects it outward. \\
Appendix C & Uses rapid shear variations and total-preserving repairs to
make the wave requirements fit. \\
\end{longtable}
}

The
\href{https://cdn.openai.com/pdf/32d9f210-8b73-45e0-91bc-82a30aef8a9a/navier-stokes.pdf\#page=13}{paper's
contents and proof-outline figures} are the source for this map. Our
examples illustrate the jobs performed by these blocks; they are not
substitutes for their proofs.

\subsection{A first reading route}\label{a-first-reading-route}

Follow these nine articles in order. Each introduces the concepts it
uses and ends by returning to the paper.

\begin{enumerate}
\def\labelenumi{\arabic{enumi}.}
\tightlist
\item
  \href{../d2_finite_energy/article.html}{How can speed rise while
  energy falls?}
\item
  \href{../d6_inner_profile/article.html}{What should happen at the
  centre of a vortex?}
\item
  \href{../d6_axis_pressure/article.html}{Why the pressure at the centre
  depends on the outside}
\item
  \href{../d3_residual/article.html}{The missing push: what a residual
  means}
\item
  \href{../convex-integration/article.html}{Convex integration: useful
  fine detail}
\item
  \href{../d4_mean_stress/article.html}{How motion with zero average can
  transport momentum}
\item
  \href{../d5_shearing_wave/article.html}{How a small wave grows, then
  fades}
\item
  \href{../d6_moment_correction/article.html}{How small local repairs
  preserve important totals}
\item
  \href{../completing-the-flow/article.html}{How the pieces become one
  complete flow}
\end{enumerate}

For a first pass, the numerical side articles are optional:
\href{../d6_global_axis/article.html}{local checks versus global
guarantees} and \href{../d7_conditioning/article.html}{what computer
digits can hide}.

The physical interpretation has its own short route:
\href{../d1_core_size/article.html}{core size versus speed},
\href{../d9_regime_map/article.html}{which small length matters}, and
\href{../d10_compressible_vortex/article.html}{what changes when density
can vary}. These ask what the model can mean for matter; they are not
additional proof steps.

\subsection{A small vocabulary to keep
nearby}\label{a-small-vocabulary-to-keep-nearby}

{\def\LTcaptype{none} % do not increment counter
\begin{longtable}[]{@{}
  >{\raggedright\arraybackslash}p{(\linewidth - 2\tabcolsep) * \real{0.5000}}
  >{\raggedright\arraybackslash}p{(\linewidth - 2\tabcolsep) * \real{0.5000}}@{}}
\toprule\noalign{}
\begin{minipage}[b]{\linewidth}\raggedright
Word
\end{minipage} & \begin{minipage}[b]{\linewidth}\raggedright
Meaning in this guide
\end{minipage} \\
\midrule\noalign{}
\endhead
\bottomrule\noalign{}
\endlastfoot
Profile & The shape of a quantity as position changes. \\
Shear & Neighbouring fluid moving at different speeds. \\
Oscillation & A repeated variation, such as a wave or a ripple. \\
Momentum & Mass times velocity; both amount and direction matter. \\
Stress & Momentum transport, or the corresponding force-per-area
description. \\
Residual & The force required by a proposed velocity and pressure;
compared with a specified force, it reveals the remaining mismatch. \\
Mean & An average over a stated group, region, or wave cycle. \\
Moment & A total in which contributions receive specified weights. \\
Smooth & Having well-behaved rates of change of every order used by the
theorem. \\
Localization & Making a field vanish outside a chosen region, with a
smooth transition. \\
\end{longtable}
}

\subsection{What understanding the paper should feel
like}\label{what-understanding-the-paper-should-feel-like}

You do not need to remember every symbol. You should be able to say what
problem each block solves, why that problem appears, and what the next
block needs from it.

The central thread is the balance between a growing velocity and a
well-behaved force. The shrinking geometry makes the energy possible;
the background establishes the shape; the ripples and pulses provide the
missing transport; the corrections and final estimates make the pieces
fit.




\end{document}
